\begin{tabularx}{\textwidth}{@{}P{0.5cm}P{1.6cm}P{3.0cm}YP{1.8cm}@{}}
\toprule
\textbf{Id} & \textbf{Tipo de dúvida} & \textbf{Pergunta (como aparece)} & \textbf{Resposta} & \textbf{Modelo de origem} \\
\midrule
H1 & Informativa; motivacional & O que é uma rotina? & ``Uma rotina é uma regra que o aplicativo segue sozinho: num horário ou numa situação, ela liga ou desliga aparelhos. Quem cria a rotina escolhe quando ela age. As rotinas ficam na aba Rotinas.'' & Esquema de signos \\
H2 & Descritiva; interpretativa & O que faz o botão de duas barras? & ``É o botão Pausar. Uma rotina pausada fica guardada, mas não age até ser retomada. Quando ela está pausada, o botão mostra um triângulo: é o Retomar.'' & Esquema de signos \\
H3 & Descritiva & O que faz o botão de lápis? & ``É o botão Editar. Ele abre a rotina para mudar o horário e os aparelhos incluídos.'' & Esquema de signos \\
H4 & Interpretativa & O que é ``Smart Schedule''? & ``É o horário em que a rotina age. Toque em $+$ ou $-$ para mudar de hora em hora, de 00h a 23h. Depois, toque em Salvar.'' & Esquema de signos; CTT \\
H5 & Investigativa; histórica & Como saber quem mudou um aparelho? & ``Abra a aba Histórico. Cada linha mostra a hora, o aparelho, o que mudou e quem mudou: `Você' ou o nome da rotina.'' \textit{[B: ``Você também pode tocar em `Por quê?' no aviso da tela Início.'']} & MoLIC (S5, S8) \\
H6 & Procedimental & Como fazer uma rotina parar de agir por um tempo? & ``Abra Rotinas, toque no nome da rotina e toque no botão de duas barras (Pausar). Para ela voltar a agir, toque no triângulo (Retomar).'' & CMN-GOMS (método PAUSAR) \\
H7 & De escolha & Pausar, editar ou excluir: qual usar? & ``Pausar: a rotina para de agir, mas fica guardada. Editar: muda o horário ou os aparelhos. Excluir: apaga a rotina; o aplicativo pede confirmação e, logo depois, ainda mostra `Desfazer' na lista de rotinas.'' & GOMS (regras de seleção); matriz de rupturas \\
H8 & Navegacional & Onde ficam as rotinas? & ``Na aba Rotinas, na barra de baixo, que aparece em todas as telas.'' & MoLIC (acesso ubíquo) \\
H9 & Procedimental (recuperação) & Como desfazer uma mudança? & ``Se você ligou ou desligou um aparelho sem querer, toque no interruptor de novo. Se pausou uma rotina, toque em Retomar. Se excluiu uma rotina, toque em Desfazer na lista de rotinas.'' \textit{[B: ``Se uma rotina mudou um aparelho, toque em Desfazer no aviso.'']} & Matriz de rupturas (RA) \\
H10 & Informativa (protocolo) & O que são Concluir e Desistir? & ``Concluir: avisa que você terminou o que foi pedido. Desistir: encerra sem terminar.'' & Protocolo \\
\bottomrule
\end{tabularx}
